\begin{statementbox}
\begin{center}
\begin{tikzpicture}[font=\small,>=Stealth,line cap=round,line join=round]
\path[use as bounding box] (-6.5,-3.55) rectangle (6.5,3.55);
\draw[black!20,rounded corners=3pt] (-6.5,-3.55) rectangle (6.5,3.55);
\node[font=\small\bfseries] at (0,3.08) {轴对称：每个点都有仍在曲线上的镜像};

\begin{scope}[shift={(-2,-1.95)},x=1.3cm,y=.85cm]
\draw[->,black!45] (-1,0)--(4.4,0) node[right] {$x$};
\draw[->,black!45] (0,-.2)--(0,4.8) node[above] {$y$};
\draw[CStatement,dashed,thick] (2,-.1)--(2,4.7) node[above right] {$x=2$};
\draw[CExplanation,very thick,domain=.15:3.85,samples=100] plot (\x,{(\x-2)^2+1});
\node[anchor=west] at (3.5,1) {$y=(x-2)^2+1$};
\draw[<->,CStatement,thick] (1,2)--(3,2);
\fill[CExplanation] (1,2) circle(2pt) node[below left=3pt] {$P$};
\fill[CExplanation] (3,2) circle(2pt) node[below right=3pt] {$P'$};
\draw[<->] (1,.45)--(2,.45) node[midway,below] {$1$};
\draw[<->] (2,.45)--(3,.45) node[midway,below] {$1$};
\end{scope}

\node[align=center] at (0,-3.05) {$P(1,2)$ 与 $P'(3,2)$ 同高、到轴等距；证明须覆盖任意点。};
\end{tikzpicture}
\end{center}


\textbf{\textcolor{CStatement}{条件}}
\begin{description}[style=nextline, leftmargin=2.8em, labelsep=0.5em, itemsep=0.35em]
  \item[\textbf{\textcolor{CStatement}{条件1（曲线与候选轴）：}}] 设曲线 $\mathcal C$ 是平面上满足 $F(x,y)=0$ 的\textbf{全部实点}，候选直线为 $\ell:Ax+By+C=0$，其中 $A^2+B^2>0$。
  \item[\textbf{\textcolor{CStatement}{条件2（对称点）：}}] 对任意点 $P(x,y)$，以 $\ell$ 为对称轴所得到的点记为 $P'(x',y')$。若是函数 $y=f(x)$，还需关注函数定义域是否一同对称。
\end{description}

\textbf{\textcolor{CStatement}{结论}}
\begin{description}[style=nextline, leftmargin=2.8em, labelsep=0.5em, itemsep=0.45em]
  \item[\textbf{\textcolor{CStatement}{结论一（根本判据）：}}]
  \textbf{直观描述：} 如图，每个点都须有仍在曲线上的镜像；只展示一对点还不能代替证明。精确写作
  \[
    \boxed{F(x,y)=0\ \Longleftrightarrow\ F(x',y')=0\quad\text{（对所有 }x,y\text{）}}.
  \]
  实际检验只须证明对任意 $P\in\mathcal C$ 有 $P'\in\mathcal C$：因为再对称一次会回到 $P$，所以反向关系自动成立。
  \item[\textbf{\textcolor{CStatement}{结论二（四类特殊轴）：}}]
  \[
  \begin{array}{c|c}
  \text{候选对称轴}&(x,y)\longmapsto(x',y')\\\hline
  x=h&(2h-x,y)\\
  y=k&(x,2k-y)\\
  y=x&(y,x)\\
  y=-x&(-y,-x)
  \end{array}
  \]
  把右栏代入 $F$：若得到与原曲线\textbf{相同的解集}，轴对称成立。尤其要区分对称轴 $y=x$ 与 $y=-x$ 的代换。
  \item[\textbf{\textcolor{CStatement}{结论三（函数图像的竖直对称轴）：}}]
  对 $y=f(x)$，直线 $x=h$ 是对称轴，当且仅当定义域关于 $h$ 对称，并且
  \[
    \boxed{f(h+t)=f(h-t)}
  \]
  对所有使两边有定义的 $t$ 成立。换言之，把自变量中心移到 $h$ 后，函数关于新原点左右相等。
  \item[\textbf{\textcolor{CStatement}{结论四（任意直线通式）：}}]
  设 $s=\dfrac{Ax+By+C}{A^2+B^2}$，则
  \[
    \boxed{x'=x-2As,\qquad y'=y-2Bs}.
  \]
  这给出任意斜率、包括竖直轴在内的统一验证方法。
  \item[\textbf{\textcolor{CStatement}{结论五（找候选轴的捷径）：}}]
  一般抛物线先找顶点和开口方向；标准椭圆、双曲线先找中心及两主轴；\textbf{真正的双曲线}两条渐近线的角平分线是两条对称轴；\textbf{一般曲线}即使有两条渐近线，也不能据此直接断定角平分线是对称轴。
\end{description}
\end{statementbox}
